\documentclass[10pt]{article}

% ---------- Document Identity (update these only) ----------
\newcommand{\DocStage}{}
\newcommand{\DocTitle}{Your Road to Tier One SOF}
\newcommand{\ProgramName}{182-Week Civilian-to-Selection Readiness Pipeline}
\newcommand{\ProgramSubtitle}{}

% ---------- Page ----------
\usepackage[legalpaper,left=0.5in,right=0.5in,top=0.6in,bottom=0.45in]{geometry}

% ---------- Typography ----------
\usepackage[T1]{fontenc}
\usepackage[utf8]{inputenc}
\usepackage{libertinus}
\usepackage{microtype}
\usepackage{parskip}
\usepackage{ragged2e}
\usepackage{textcomp}
\AtBeginDocument{\color{black}}
\AtBeginDocument{\pagecolor{white}}
\AtBeginDocument{\justifying}

% ---------- Landscape pages ----------
\usepackage{pdflscape}

% ---------- Color ----------
\usepackage[table]{xcolor}
\definecolor{StageBlue}{HTML}{1F4E79}
\definecolor{StageBlueDark}{HTML}{163A5A}
\definecolor{LightGray}{HTML}{F2F4F7}
\definecolor{MidGray}{HTML}{D9DEE5}
\definecolor{RuleGray}{HTML}{C7CED8}
\definecolor{HeaderInk}{HTML}{000000}
\definecolor{Ink}{HTML}{000000}

% ---------- Utility ----------
\usepackage{enumitem}
\sloppy
\setlength{\emergencystretch}{2em}

% ---------- Boxes / UI ----------
\usepackage[most]{tcolorbox}

\tcbset{
  charterTitle/.style={
    enhanced,
    colback=StageBlue!4,
    colframe=StageBlue,
    boxrule=1.25pt,
    arc=3pt,
    left=16pt,right=16pt,top=14pt,bottom=14pt,
    borderline north={3.2pt}{0pt}{StageBlue},
    borderline west={4pt}{0pt}{StageBlue},
    borderline south={1.25pt}{0pt}{StageBlue},
    drop shadow={black!25!white},
  },
  charterRule/.style={
    enhanced,
    colback=LightGray,
    colframe=RuleGray,
    boxrule=0.85pt,
    arc=3pt,
    left=12pt,right=12pt,top=10pt,bottom=10pt,
    borderline west={3pt}{0pt}{StageBlue},
  },
  charterBadge/.style={
    enhanced,
    colback=StageBlue,
    coltext=white,
    colframe=StageBlueDark,
    boxrule=0.6pt,
    arc=2pt,
    left=6pt,right=6pt,top=3pt,bottom=3pt,
  }
}
\newtcbox{\Badge}{nobeforeafter,charterBadge}

% ---------- Lists ----------
\setlist[itemize]{leftmargin=1.5em, itemsep=0.25em, topsep=0.25em}
\setlist[enumerate]{leftmargin=1.8em, itemsep=0.25em, topsep=0.25em}

% ---------- TOC ----------
\usepackage{tocloft}
\setcounter{tocdepth}{3}
\setlength{\cftbeforesecskip}{3pt}
\setlength{\cftbeforesubsecskip}{2pt}
\setlength{\cftbeforesubsubsecskip}{0pt}
\renewcommand{\cftsecfont}{\bfseries}
\renewcommand{\cftsecpagefont}{\bfseries}
\renewcommand{\cftsubsecfont}{\normalfont}
\renewcommand{\cftsubsecpagefont}{\normalfont}
\renewcommand{\cftsubsubsecfont}{\normalfont}
\renewcommand{\cftsubsubsecpagefont}{\normalfont}
\setlength{\cftsecindent}{0pt}
\setlength{\cftsubsecindent}{1.25em}
\setlength{\cftsubsubsecindent}{2.8em}
\setlength{\cftsecnumwidth}{2.1em}
\setlength{\cftsubsecnumwidth}{2.8em}
\setlength{\cftsubsubsecnumwidth}{3.6em}

% Remove the built-in TOC heading so only the custom heading is shown
\makeatletter
\renewcommand{\tableofcontents}{\@starttoc{toc}}
\makeatother

% ---------- Header/Footer: page number only ----------
\usepackage{fancyhdr}
\pagestyle{fancy}
\fancyhf{}
\renewcommand{\headrulewidth}{0pt}
\renewcommand{\footrulewidth}{0pt}
\fancyfoot[C]{\thepage}

% ---------- Section formatting ----------
\usepackage{titlesec}

\titleformat{\section}
  {\Large\bfseries\color{Ink}}
  {\thesection}{0.6em}{}
  [\vspace{0.25em}\textcolor{StageBlue}{\rule{\linewidth}{0.8pt}}]

\titleformat{\subsection}
  {\large\bfseries\color{Ink}}
  {\thesubsection}{0.6em}{}

\titleformat{\subsubsection}
  {\normalsize\bfseries\color{Ink}}
  {\thesubsubsection}{0.6em}{}

\titlespacing*{\section}{0pt}{0.95em}{0.35em}
\titlespacing*{\subsection}{0pt}{0.65em}{0.25em}
\titlespacing*{\subsubsection}{0pt}{0.55em}{0.2em}

% ---------- Table engine ----------
\usepackage{tabularray}
\UseTblrLibrary{booktabs}

% ---------- tabularray: GLOBAL table treatment ----------
\SetTblrInner{
  cells = {fg=black, bg=white, valign=t},
}

% ---------- Header helpers ----------
\newcommand{\Hdr}[1]{\shortstack[c]{\strut #1\strut}}

% ---------- Lines ----------
\newcommand{\TitleLine}{\textcolor{StageBlue}{\rule{0.98\linewidth}{0.95pt}}}
\newcommand{\SubTitleLine}{\textcolor{RuleGray}{\rule{0.98\linewidth}{0.65pt}}}

% ---------- Continued on next page ----------
\newcommand{\ContinuedOnNextPageInline}{%
  \par
  \vfill
  \noindent\textcolor{RuleGray}{\rule{\linewidth}{1pt}}\vspace{-2pt}
  \noindent\hfill\textit{Continued on next page}\par\vspace{-10pt}
  \noindent\textcolor{RuleGray}{\rule{\linewidth}{1pt}}
  \clearpage
}

% ---------- PDF hyperlinks ----------
\usepackage[hidelinks]{hyperref}
\hypersetup{
  colorlinks=false,
  pdfborder={0 0 0},
  linktoc=all,
  pdftitle={\DocTitle},
  pdfauthor={\ProgramName}
}

% =========================================================
\begin{document}

\section*{Stage 5.B — Detailed Phase Card — Phase 04 Card Header}
\phantomsection
\addcontentsline{toc}{subsection}{Phase 04 Card Header}

\begin{itemize}
\item Phase \textbf{ID}: Phase 04
\item Phase Name: Strength-Endurance and Work Capacity Development
\item Duration weeks: 12
\item Version: v1.0
\item Governing Status: Deployable
\end{itemize}

\subsection*{1. Phase Mission}
\phantomsection


Phase 04 establishes strength-endurance and work-capacity development while introducing the first default gate-relevant load-carriage time trial construct (\textbf{RUCK} 6-mile, 35 lb dry) under strict admissibility and comparability controls. The phase must prove that ruck evidence can be produced without foot breakdown, without gait drift, and without validity contamination, while sustaining \textbf{RUN} and \textbf{CAL} verification continuity under protected windows.

Mission boundaries (binding):

\begin{itemize}
\item Phase 04 does not authorize unsafe load escalation or informal ruck testing outside protected windows to “see where we’re at.”
\item Phase 04 does not treat ruck evidence as interchangeable with \textbf{RUN} evidence and does not permit “equivalent conditioning” substitutions to claim ruck readiness.
\item Phase 04 does not permit “stack and salvage” behavior inside protected windows; when caps/spacing cannot be preserved, the correct response is reslot.
\item Phase 04 does not weaken Stage 0.5 safety posture or Stage 1.F admissibility rules.
\end{itemize}

\subsection*{2. Phase Purpose}
\phantomsection


Phase 04 exists to solve a specific transition problem in the macro pipeline: moving from Phase 03’s endurance-conversion emphasis into Phase 05’s extended stamina emphasis requires a foot, route, and load interface that is stable enough to support load carriage without predictable breakdown. Phase 04 therefore formalizes:

A) Ruck admissibility as a capability, not a preference

\begin{itemize}
\item The ruck construct requires:
  \begin{itemize}
  \item stable footwear interface,
  \item stable footcare protocol execution,
  \item dry-load verification posture,
  \item stable route or treadmill identity and logging,
  \item stable gait mechanics.
  \end{itemize}
\end{itemize}

If these are not present, Phase 04 cannot be considered gate-ready, regardless of other domains.

B) Strength-endurance and work capacity as phase intent

\begin{itemize}
\item Work-capacity improvements must occur without converting Phase 04 into breathless, unscheduled “benchmark culture.”
\item The phase increases the ability to sustain output under controlled effort while preserving recovery enough to produce valid test evidence.
\end{itemize}

C) Preservation of \textbf{RUN} and \textbf{CAL} evidence streams

\begin{itemize}
\item \textbf{RUN} remains a separate, gate-relevant domain and must not be deprioritized simply because ruck is introduced.
\item \textbf{CAL} remains a standardized battery with artifact and validity posture preserved; it continues to provide a stable readiness signal across phases.
\end{itemize}

Not the purpose (explicit drift prevention):

\begin{itemize}
\item Not to introduce new tests, simulations, mixed-modal benchmarks, or interval benchmarks beyond Stage 2 instruments.
\item Not to shift gate-grade tests into build weeks.
\item Not to pursue maximal strength testing; \textbf{STR} remains supporting and submax by posture.
\end{itemize}

\subsection*{3. Entry Criteria (Gates to Start)}
\phantomsection


Entry criteria are objective and auditable. If unmet, the default posture is \textbf{HOLD} at phase transition with reslot and stabilization.

A) Prior phase outcome and evidence integrity

\begin{itemize}
\item Phase 03 end gate outcome is \textbf{ADVANCE}, recorded under Stage 1.F.
\item The phase transition is not permitted if Phase 03 evidence is contaminated:
  \begin{itemize}
  \item unresolved stop-rule events,
  \item missing minimum dataset,
  \item or gate decisions made on inadmissible evidence.
  \end{itemize}
\end{itemize}

B) Decision window compliance prerequisites (Stage 1.F default posture)

For the two most recent complete training weeks preceding Phase 04 Week 1:

\begin{itemize}
\item Admissible sessions: at least 5 of 6 admissible sessions per week on average.
\item \textbf{MODIFIED} admissibility cap: outside medical issues, ≤1 \textbf{MODIFIED} session per week admissible.
\item Disqualifier: any \textbf{INVALID} session in the decision window disqualifies \textbf{ADVANCE} posture and forces \textbf{HOLD}.
\end{itemize}

C) Durability prerequisites (Phase 04-specific because ruck becomes gate-relevant)

\begin{itemize}
\item \textbf{MET-2B-036} (Gait-Altering Pain Flag): No across the decision window.
\item \textbf{MET-2B-037} (Foot Hotspot or Blister Flag): no active blister and no active escalation trend that would predictably compromise ruck execution.
\item Pain posture: no unresolved mechanics-altering pain trends that would predictably force substitution.
\end{itemize}

D) Ruck feasibility prerequisites (must exist before the first ruck protected window)

\begin{itemize}
\item Load verification method is feasible and loggable (dry load only; water logged separately) per Stage 2 \textbf{C12} posture.
\item A measurable ruck environment exists:
  \begin{itemize}
  \item stable route identity, or
  \item treadmill alternative with identity/settings logging posture.
  \end{itemize}
\item Pack configuration posture can be held stable enough to support comparability (category-level description permitted; no product specifics required).
\end{itemize}

Default posture if any of the above is not satisfied:

\begin{itemize}
\item \textbf{HOLD} at phase transition, preserve cadence via conservative substitutions, and reslot the first ruck test attempt until prerequisites are validated.
\end{itemize}

\subsection*{4. Operating Constraints}
\phantomsection


These constraints apply to every week of Phase 04 and directly constrain admissibility and gate decisions.

A) Cadence and session structure (non-negotiable)

\begin{itemize}
\item 6 sessions per week is mandatory.
\item No two-a-days. No make-up doubling. Missed sessions remain missed and are logged accurately.
\item Required session elements per Stage 1.F in correct order:
  \begin{itemize}
  \item Safety self-check first
  \item Warm-up
  \item Mobility
  \item Main work
  \item Cardio
  \item Cooldown
  \item Recovery notes and post-session notes
  \end{itemize}
\end{itemize}

B) Cardio minimum and continuity (enforced in body text)

\begin{itemize}
\item Cardio occurs in every training session.
\item Cardio minimum standard is minimum 30 minutes continuous per session.
\item If Cardio is under the minimum or not continuous, the session is \textbf{INVALID} and provides no gate credit.
\end{itemize}

C) Baseline daily activity

\begin{itemize}
\item 10,000 steps per day is a global baseline activity requirement and must be logged.
\item Steps do not count as a session and do not replace session requirements.
\end{itemize}

D) Evidence posture

\begin{itemize}
\item Gate decisions use admissible evidence only:
  \begin{itemize}
  \item sessions must be admissible under Stage 1.F (\textbf{VALID}; \textbf{MODIFIED} admissible only within cap rules),
  \item tests must be \textbf{VALID} to be admissible,
  \item \textbf{INVALID} tests provide no gate credit and require reslot.
  \end{itemize}
\end{itemize}

E) Phase 04 high-consequence governance emphasis

\begin{itemize}
\item \textbf{RUCK} gate construct is treated as high-cost evidence:
  \begin{itemize}
  \item it is not executed opportunistically in build weeks as a quasi-test,
  \item it is executed only in protected windows,
  \item it is aborted/reslotted when foot integrity or environment validity cannot be preserved.
  \end{itemize}
\end{itemize}

F) Water governance

\begin{itemize}
\item No underwater or breath-hold exposure is scheduled in Phase 04; underwater is not a Phase 04 gate domain.
\item Any water governance breach is a safety event and contaminates evidence for gate purposes.
\end{itemize}

\subsection*{5. Primary Adaptations and Physiological Focus}
\phantomsection


Phase 04 adaptation goals are expressed as governed outcomes (no programming detail):

\begin{itemize}
\item Strength-endurance expansion: improved ability to sustain controlled work output without form breakdown and without recovery collapse.
\item Work-capacity development: improved tolerance for multi-domain workload while maintaining evidence integrity and stable session completion.
\item Load-carriage tolerance initiation and stabilization:
  \begin{itemize}
  \item proof that the foot/footwear system can sustain load carriage exposure without breakdown,
  \item proof that dry-load verification and logging posture are executable and stable,
  \item proof that gait mechanics remain stable under declared load.
  \end{itemize}
\item \textbf{RUN} continuity under controlled fatigue:
  \begin{itemize}
  \item maintain and improve \textbf{RUN} evidence without allowing ruck fatigue to contaminate \textbf{RUN} test windows.
  \end{itemize}
\item \textbf{CAL} standardization continuity:
  \begin{itemize}
  \item maintain strict standards and artifact validity posture for gate-grade \textbf{CAL} evidence.
  \end{itemize}
\item Durability and recovery stability as gate constraints:
  \begin{itemize}
  \item stable pain trends, stable gait mechanics, controlled foot hotspot incidence,
  \item sleep/fatigue posture stable enough to avoid predictable invalidity in protected windows.
  \end{itemize}
\item \textbf{BODYCOMP} trend continuity as risk-control context:
  \begin{itemize}
  \item trend directionality supports feasibility and reduces injury risk; it does not substitute for performance evidence.
  \end{itemize}
\end{itemize}

\subsection*{6. Primary Modalities (Strength, Conditioning, Ruck, Swim, and Other)}
\phantomsection


Domain emphasis and governance posture:

A) Primary emphasized domains

\begin{itemize}
\item \textbf{RUCK}: gate-relevant construct introduced; strict dry-load verification posture and environment identity posture.
\item \textbf{RUN}: maintained and tested under protected windows; must remain separate from ruck constructs.
\item \textbf{CAL}: standardized battery maintained; must remain artifact-valid for gate use.
\item \textbf{DURABILITY}: constraint signals determine eligibility for gate-grade attempts.
\item \textbf{RECOVERY}: readiness markers used to decide whether a protected week remains test-capable.
\end{itemize}

B) Supporting domains

\begin{itemize}
\item \textbf{STR}: supporting only; submax posture; must not compromise ruck/run validity through fatigue spillover.
\item \textbf{BODYCOMP}: continuous trend capture for feasibility and risk control.
\end{itemize}

C) De-emphasized domains (Phase 04 default posture)

\begin{itemize}
\item \textbf{SWIM}: not required for Phase 04 gating by default; permitted only if pool access and unit verification posture are stable and adding it does not violate cap and spacing.
\item \textbf{AER}: not used as \textbf{RUN} equivalency; may be used as low-impact Cardio modality to preserve cadence and safety under foot breakdown or environment hazards.
\end{itemize}

\subsection*{7. Weekly Structure}
\phantomsection


Intent-level weekly structure maintains phase contract without becoming a program.

\clearpage
\begin{landscape}

\begingroup
\scriptsize
\setlength{\tabcolsep}{2pt}
\renewcommand{\arraystretch}{1.12}

\begin{longtblr}{
  width=\linewidth,
  rowhead=1,
  colspec = {
    Q[l,wd=0.70in]
    X[l,1.75]
    X[l,1.25]
    X[l,2.75]
    X[l,2.75]
  },
  hlines,
  vlines,
  cells = {hsep=6pt, vsep=6pt, halign=l, valign=t},
  cell{1-Z}{1-5} = {fg=black},
  row{odd} = {bg=white},
  row{even} = {bg=LightGray},
  row{1} = {bg=StageBlue!15, fg=black, font=\bfseries, halign=l, valign=t},
}
\Hdr{Session \#} &
\Hdr{Primary Focus} &
\Hdr{Main Work Domains} &
\Hdr{Cardio Modality/Intent} &
\Hdr{Notes/Risk Controls} \\

1 &
\textbf{RUN} continuity + controlled effort posture &
\textbf{RUN}; \textbf{DURABILITY}; \textbf{RECOVERY} &
Modality: run or treadmill run where safe; Minutes: 30–60; RPE + talk-test: RPE 3–5, full-to-short sentences; Continuity: continuous planned &
Not a time trial day. Surface and footwear stability enforced. If \textbf{MET-2B-036} becomes Yes or \textbf{MET-2B-037} escalates, substitute to \textbf{AER} modality and document; preserve Cardio continuity and safety. \\

2 &
Strength-endurance patterning + mobility bias &
\textbf{STR}; \textbf{DURABILITY} &
Modality: low-impact steady (walk/bike/elliptical); Minutes: 30–45; RPE + talk-test: RPE 3–4 full sentences; Continuity: continuous planned &
\textbf{STR} is supporting only. Avoid fatigue spikes that contaminate ruck readiness. Terminate or regress on mechanics change. \\

3 &
\textbf{RUCK} exposure governance day (non-test) &
\textbf{RUCK}; \textbf{DURABILITY}; \textbf{RECOVERY} &
Modality: ruck exposure only if foot integrity stable; otherwise substitute; Minutes: 30–60; RPE + talk-test: RPE 3–5; Continuity: continuous planned &
High-risk day. Foot hotspot monitoring is mandatory. Any mechanics change ends the exposure; do not convert to “finish anyway.” Protect skin integrity; substitution is preferred over omission. \\

4 &
\textbf{CAL} integrity maintenance + trunk stability &
\textbf{CAL}; \textbf{DURABILITY} &
Modality: low-impact steady; Minutes: 30–45; RPE + talk-test: RPE 2–4 full sentences; Continuity: continuous planned &
Maintain standards and technique without maximal attempts. Preserve shoulders/elbows/wrists. Prepare artifact workflow for protected windows without forcing maximal efforts mid-week. \\

5 &
Aerobic throughput buffer + recovery protection &
\textbf{RECOVERY}; \textbf{DURABILITY} &
Modality: lowest-risk continuous Cardio; Minutes: 30–45; RPE + talk-test: RPE 2–4 full sentences; Continuity: continuous planned &
Buffer day. Used to protect the fatigue budget and preserve compliance. When readiness is degraded, treat as Minimum Viable Session posture while preserving structure and logging. \\

6 &
Consolidation + pre-window stabilization &
\textbf{RUN} or \textbf{RUCK} (light); \textbf{RECOVERY} &
Modality: step-based or indoor substitute; Minutes: 30–45; RPE + talk-test: RPE 2–4 full sentences; Continuity: continuous planned &
Stabilization day before protected windows. Confirm environment/tool identity readiness, load verification readiness for ruck window, and complete logging readiness. No novelty in footwear or tools. \\

\end{longtblr}

\endgroup

\end{landscape}
\clearpage

\subsection*{8. Progression Rules (Volume and Intensity Caps; Regression Triggers)}
\phantomsection


A) Progression doctrine (governed; conservative)

\begin{itemize}
\item Phase 04 progression is stability-driven and must preserve admissibility:
  \begin{itemize}
  \item stable mechanics,
  \item stable foot integrity,
  \item stable recovery posture,
  \item stable environment or tool identity posture.
  \end{itemize}
\item One progression lever at a time:
  \begin{itemize}
  \item do not escalate ruck exposure dose and increase \textbf{RUN} test intensity posture in the same week,
  \item do not increase total weekly Cardio minutes and add higher-cost load exposures simultaneously without demonstrated stability.
  \end{itemize}
\end{itemize}

B) \textbf{RUCK} progression safeguards (Phase 04 specific)

\begin{itemize}
\item Ruck exposure escalation is prohibited when:
  \begin{itemize}
  \item gait-altering pain flag is Yes,
  \item active blister or hotspot escalation trend exists,
  \item load verification cannot be executed and logged,
  \item environment is unsafe or measurability cannot be preserved.
  \end{itemize}
\item If a ruck exposure creates hotspot escalation:
  \begin{itemize}
  \item treat as a durability constraint,
  \item downshift next exposure and consider substituting to low-friction modalities,
  \item reslot any planned ruck time trial until stability is restored.
  \end{itemize}
\end{itemize}

C) \textbf{RUN} progression safeguards (interaction with ruck)

\begin{itemize}
\item \textbf{RUN} verification must not be scheduled or executed under fatigue spillover from ruck:
  \begin{itemize}
  \item preserve spacing in protected windows,
  \item avoid high-cost ruck exposures immediately before \textbf{RUN} test days in protected windows (staging intent without day-level schedule).
  \end{itemize}
\end{itemize}

D) Regression triggers (hard controls)

Same-day regression triggers (terminate provoking exposure):

\begin{itemize}
\item pain that alters mechanics,
\item dizziness/lightheadedness,
\item unusual/disproportionate breathlessness relative to planned effort,
\item instability/swelling,
\item stop-rule symptoms per Stage 0.5 / Stage 1.F,
\item environment hazards that create slip/fall risk or unsafe exertion risk.
\end{itemize}

Trend-based triggers (convert protected window to deload-only/check-in-only; reslot maximal tests):

\begin{itemize}
\item sustained sleep debt trend with rising session RPE at stable workload,
\item rising soreness/fatigue trends that predict invalid tests,
\item repeated \textbf{MODIFIED}/\textbf{INVALID} sessions in the decision window,
\item tool or environment drift not stabilized.
\end{itemize}

\subsection*{9. Deload and Test Placement (Mid-Phase and End-of-Phase)}
\phantomsection


Phase 04 is 12 weeks. Protected windows occur per Stage 3 doctrine.

Protected deload/test windows (week-in-phase):

\begin{itemize}
\item Week 4: Deload+Test mid-phase window (primary: ruck proof event by default)
\item Week 8: Deload+Test mid-phase window (primary: \textbf{RUN} proof event by default)
\item Week 12: Deload+Test end-phase exit gate window (staged battery)
\end{itemize}

Major event posture inside protected windows:

\begin{itemize}
\item Ideal cap: 2 major test events per week.
\item Hard cap: 3 major test events per week.
\item Spacing intent:
  \begin{itemize}
  \item any two major test events must not be scheduled within the same 72-hour window (start-to-start),
  \item maximal \textbf{RUN} and maximal \textbf{RUCK} must not be scheduled within the same 96-hour window (start-to-start).
  \end{itemize}
\end{itemize}

When caps/spacing cannot be preserved:

\begin{itemize}
\item Do not compress.
\item Complete the decisive items and reslot the remainder to the next protected window per Stage 3.E (including phase extension under change control if needed).
\end{itemize}

\subsection*{10. Test Battery (IDs and Labels Only)}
\phantomsection


Gate-grade tests are referenced by Stage 2 IDs only.

\subsubsection*{Mid-Phase Tests}
\phantomsection

Week 4 mid-phase protected window (primary ruck proof event posture):

\begin{itemize}
\item \textbf{RUCK}: \textbf{C12} / \textbf{MET-2B-019} 6-mile ruck time trial (35 lb dry)
\item \textbf{CAL} (optional within cap and spacing feasibility): \textbf{C4} / \textbf{MET-2B-010}; \textbf{C5} / \textbf{MET-2B-011}; \textbf{C6} / \textbf{MET-2B-012}; \textbf{C7} / \textbf{MET-2B-013}
\end{itemize}

Week 8 mid-phase protected window (primary \textbf{RUN} proof event posture):

\begin{itemize}
\item \textbf{RUN}: \textbf{C10} / \textbf{MET-2B-016} 3-mile run time trial
\item \textbf{CAL} (optional within cap and spacing feasibility): \textbf{C4} / \textbf{MET-2B-010}; \textbf{C5} / \textbf{MET-2B-011}; \textbf{C6} / \textbf{MET-2B-012}; \textbf{C7} / \textbf{MET-2B-013}
\end{itemize}

\subsubsection*{End-Phase Tests (Week 12 exit gate)}
\phantomsection

Staged end-phase battery (subject to cap/spacing and prioritization posture):

\begin{itemize}
\item \textbf{RUN}: \textbf{C9} / \textbf{MET-2B-015} 2-mile run time trial
\item \textbf{RUCK}: \textbf{C12} / \textbf{MET-2B-019} 6-mile ruck time trial (35 lb dry)
\item \textbf{CAL} battery: \textbf{C4} / \textbf{MET-2B-010}; \textbf{C5} / \textbf{MET-2B-011}; \textbf{C6} / \textbf{MET-2B-012}; \textbf{C7} / \textbf{MET-2B-013}
\end{itemize}

\subsubsection*{Test Battery Table}
\phantomsection

\clearpage
\begin{landscape}

\begingroup
\scriptsize
\setlength{\tabcolsep}{2pt}
\renewcommand{\arraystretch}{1.12}

\begin{longtblr}{
  width=\linewidth,
  rowhead=1,
  colspec = {
    Q[l,wd=0.90in]
    X[l,2.35]
    Q[l,wd=0.85in]
    X[l,2.95]
    X[l,1.45]
  },
  hlines,
  vlines,
  cells = {hsep=6pt, vsep=6pt, halign=l, valign=t},
  cell{1-Z}{1-5} = {fg=black},
  row{odd} = {bg=white},
  row{even} = {bg=LightGray},
  row{1} = {bg=StageBlue!15, fg=black, font=\bfseries, halign=l, valign=t},
}
\Hdr{Test Timing} &
\Hdr{Test \textbf{ID}/Name} &
\Hdr{Domain} &
\Hdr{Validity Conditions} &
\Hdr{Pass Threshold} \\

Week 4 &
\textbf{C12} / \textbf{MET-2B-019} 6-mile ruck time trial (35 lb dry) &
\textbf{RUCK} &
Dry load verified and logged; water logged separately; route/treadmill identity logged; warm-up logged; safety screen passed; no stop-rule violations; Evidence Ref recorded; required fields complete &
Tier per Stage 2 (mid-phase check-in posture) \\

Week 4 &
\textbf{CAL} battery: \textbf{C4} / \textbf{MET-2B-010}; \textbf{C5} / \textbf{MET-2B-011}; \textbf{C6} / \textbf{MET-2B-012}; \textbf{C7} / \textbf{MET-2B-013} &
\textbf{CAL} &
Valid only when artifact requirements and required fields are complete; Evidence Ref recorded; missing required items defaults \textbf{INVALID} for gate use &
Tier per Stage 2 (mid-phase check-in posture) \\

Week 8 &
\textbf{C10} / \textbf{MET-2B-016} 3-mile run time trial &
\textbf{RUN} &
Route or treadmill identity logged; treadmill settings logged when used; warm-up logged; safety screen passed; Evidence Ref recorded; required fields complete; reslot if unsafe environment &
Tier per Stage 2 (mid-phase check-in posture) \\

Week 8 &
\textbf{CAL} battery: \textbf{C4} / \textbf{MET-2B-010}; \textbf{C5} / \textbf{MET-2B-011}; \textbf{C6} / \textbf{MET-2B-012}; \textbf{C7} / \textbf{MET-2B-013} &
\textbf{CAL} &
Artifact and required fields complete; Evidence Ref recorded &
Tier per Stage 2 (mid-phase check-in posture) \\

Week 12 &
\textbf{C9} / \textbf{MET-2B-015} 2-mile run time trial &
\textbf{RUN} &
Gate-grade: must be \textbf{VALID}; route or treadmill identity logged; required settings logged when treadmill is used; warm-up logged; Evidence Ref recorded; no stop-rule violations &
\textbf{INTERMEDIATE} tier per Stage 2 (with replication posture where required) \\

Week 12 &
\textbf{C12} / \textbf{MET-2B-019} 6-mile ruck time trial (35 lb dry) &
\textbf{RUCK} &
Gate-grade: must be \textbf{VALID}; dry load verified; environment identity logged; footwear recorded; warm-up logged; Evidence Ref recorded; no stop-rule violations &
\textbf{INTERMEDIATE} tier per Stage 2 (with replication posture where required) \\

Week 12 &
\textbf{CAL} battery: \textbf{C4} / \textbf{MET-2B-010}; \textbf{C5} / \textbf{MET-2B-011}; \textbf{C6} / \textbf{MET-2B-012}; \textbf{C7} / \textbf{MET-2B-013} &
\textbf{CAL} &
Gate-grade: must be \textbf{VALID}; artifact posture satisfied; Evidence Ref recorded &
\textbf{INTERMEDIATE} tier per Stage 2 (with replication posture where required) \\

\end{longtblr}

\endgroup

\end{landscape}
\clearpage

\subsection*{11. Exit Standards (Objective Pass and Fail)}
\phantomsection


Phase 04 exit gate occurs in Week 12 protected window.

A) Eligibility to decide (Stage 1.F decision window posture)

\begin{itemize}
\item Decision window: the two most recent complete training weeks preceding Week 12.
\item Minimum admissible sessions: at least 5 of 6 admissible sessions per week on average.
\item \textbf{MODIFIED} cap: outside medical issues, ≤1 \textbf{MODIFIED} session per week admissible.
\item Disqualifier: any \textbf{INVALID} session in the decision window disqualifies \textbf{ADVANCE} and forces \textbf{HOLD}.
\end{itemize}

B) Required gate-grade tests (Week 12; \textbf{VALID} only)

\textbf{ADVANCE} to Phase 05 requires all of the following:

\begin{itemize}
\item \textbf{RUN} gate item:
  \begin{itemize}
  \item \textbf{C9} / \textbf{MET-2B-015} is \textbf{VALID} and meets \textbf{INTERMEDIATE} tier posture per Stage 2.
  \end{itemize}
\item \textbf{RUCK} gate item:
  \begin{itemize}
  \item \textbf{C12} / \textbf{MET-2B-019} is \textbf{VALID} and meets \textbf{INTERMEDIATE} tier posture per Stage 2, with dry-load verification and stable environment identity.
  \end{itemize}
\item \textbf{CAL} gate items:
  \begin{itemize}
  \item \textbf{C4} / \textbf{MET-2B-010}; \textbf{C5} / \textbf{MET-2B-011}; \textbf{C6} / \textbf{MET-2B-012}; \textbf{C7} / \textbf{MET-2B-013} are \textbf{VALID} and meet \textbf{INTERMEDIATE} tier posture per Stage 2 under required validity posture.
  \end{itemize}
\item Durability disqualifier posture:
  \begin{itemize}
  \item \textbf{MET-2B-036} must be No in the decision window and at gate testing,
  \item no active blister at time of gate testing,
  \item no active breakdown trend that predicts gait drift.
  \end{itemize}
\end{itemize}

C) Replication posture for tier use

\begin{itemize}
\item Where Stage 2 requires replication for tier classification used in gate decisions, apply the replication rule (2 of last 3 \textbf{VALID} events in the decision window for that metric). If replication is not met, default posture is \textbf{HOLD} for that metric until replication exists.
\end{itemize}

D) Failure disposition (mandatory)

If any required gate item is missing or \textbf{INVALID}, or if decision-window disqualifiers are present:

\begin{itemize}
\item Default posture is \textbf{HOLD}.
\item Corrective action is reslot to the next protected window per Stage 3.E (phase extension under change control if needed).
\item Do not advance on partial gate evidence.
\end{itemize}

\subsection*{12. Risk Controls and Stop Rules (Phase-Specific Overlays)}
\phantomsection


Stage 0.5 and Stage 1.F remain controlling authorities. Phase 04 overlays apply them specifically to ruck integration.

\textbf{RUCK}-specific risk controls:

\begin{itemize}
\item Treat ruck as a high-consequence exposure:
  \begin{itemize}
  \item eligibility depends on foot integrity and gait stability,
  \item load verification and environment identity are validity prerequisites,
  \item any mechanics-altering pain or hotspot escalation triggers immediate regression and reslot posture.
  \end{itemize}
\item Dry load governance is mandatory for \textbf{C12}:
  \begin{itemize}
  \item dry load is the official load; carried water is logged separately,
  \item any unverified load or load mismatch is \textbf{INVALID} for gate use.
  \end{itemize}
\end{itemize}

\textbf{RUN}-to-\textbf{RUCK} interaction controls:

\begin{itemize}
\item Enforce spacing intent in protected windows:
  \begin{itemize}
  \item do not schedule maximal ruck and maximal run within 96 hours start-to-start,
  \item if spacing cannot be preserved, reslot the lower-priority event and preserve the decisive gate item for the phase.
  \end{itemize}
\end{itemize}

\textbf{CAL} controls:

\begin{itemize}
\item \textbf{CAL} battery is treated as a major event.
\item \textbf{CAL} gate-grade results require full validity posture (including artifact requirements where applicable); missing requirements = \textbf{INVALID} for gate use.
\end{itemize}

Environment controls:

\begin{itemize}
\item Grand Forks hazards are treated as expected constraints:
  \begin{itemize}
  \item unsafe footing, poor visibility, lightning, unsafe heat, poor air quality, or facility disruptions invalidate gate-grade attempts.
  \item indoor substitutions are used only when Stage 2 validity/logging preserves comparability; otherwise reslot.
  \end{itemize}
\end{itemize}

Stop posture reminders:

\begin{itemize}
\item \textbf{STOP} \textbf{NOW} triggers terminate the attempt; tests are tagged \textbf{INVALID} with the appropriate \textbf{TEST-RC} and reslotted.
\end{itemize}

\clearpage
\begin{landscape}

\subsection*{13. Required Capabilities and Dependencies}
\phantomsection


Dependencies are expressed as capability categories (Stage 7) plus validated prerequisites (Stage 1.C IDs). Fail actions must default conservative.

\begingroup
\scriptsize
\setlength{\tabcolsep}{2pt}
\renewcommand{\arraystretch}{1.12}

\begin{longtblr}{
  width=\linewidth,
  rowhead=1,
  colspec = {
    Q[l,wd=0.95in]
    X[l,2.60]
    Q[l,wd=1.35in]
    Q[l,wd=1.15in]
    X[l,3.10]
  },
  hlines,
  vlines,
  cells = {hsep=6pt, vsep=6pt, halign=l, valign=t},
  cell{1-Z}{1-5} = {fg=black},
  row{odd} = {bg=white},
  row{even} = {bg=LightGray},
  row{1} = {bg=StageBlue!15, fg=black, font=\bfseries, halign=l, valign=t},
}
\Hdr{Dependency \textbf{ID}} &
\Hdr{Required Capability Category and Tier} &
\Hdr{Due-By week-in-phase} &
\Hdr{Owner} &
\Hdr{Fail Action} \\

\textbf{DEP-1C-007} &
7.11 Footwear Systems and Foot Care — Tier 2 stability for load carriage (ruck interface stable; no novelty near test windows) &
Week 1 and maintained through Week 12 &
Equipment Lead &
If instability or hotspots/blisters trend upward, do not attempt ruck \textbf{TT}; preserve Cardio via low-friction substitutes; reslot ruck \textbf{TT}; default \textbf{HOLD} if ruck evidence cannot be produced validly. \\

\textbf{DEP-1C-008} &
7.11 Footwear Systems and Foot Care — Tier 1 hotspot protocol execution and drying routine &
Week 1 &
Equipment Lead &
If hotspot protocol cannot be executed reliably, downshift exposure; ruck \textbf{TT} is not attempted; reslot; default \textbf{HOLD} for gate posture dependent on ruck admissibility. \\

\textbf{DEP-1C-006} &
7.18 Testing, Timing, Measurement, and Course-Setup Tools — Tier 1 route/treadmill identity and measurability posture (\textbf{ROUTE}-2E-### / \textbf{TM}-2E-###) &
Week 1 &
Coach Lead &
If identity/measurability cannot be logged, \textbf{RUN} and ruck tests are \textbf{INVALID} for gate use; reslot; do not advance on compromised measurement integrity. \\

\textbf{DEP-1C-014} &
7.08 Conditioning, Endurance, and Aerobic Training Tools — Tier 1 safe Cardio access + indoor substitute posture under environment volatility &
Week 1 &
Coach Lead &
If safe Cardio access cannot be maintained, sessions become inadmissible; default \textbf{HOLD} until feasibility restored; preserve cadence via the safest option available or terminate when unsafe. \\

\textbf{DEP-1C-019} &
7.14 Environmental Apparel — Tier 1 environment go/no-go and substitution discipline &
Week 2 &
Trainee &
If environment hazards prevent safe execution and indoor substitution cannot preserve comparability for gate tests, reslot; do not force unsafe attempts due to schedule pressure. \\

\textbf{DEP-1C-021} &
7.16 Operational Self-Management Systems — Tier 1 same-day logging compliance (minimum dataset completeness) &
Week 1 &
Trainee &
Missing logs make evidence inadmissible; default \textbf{HOLD} until compliance restored; repeat compliance proof; do not attempt gate-grade tests without evidence integrity. \\

\textbf{DEP-1C-016} &
7.10 Logistics, Maintenance, Storage, and Sustainment Equipment — Tier 1 drying and sustainment posture &
Week 2 &
Trainee &
If drying/sustainment fails and foot breakdown recurs, downshift; reslot; treat as durability constraint that blocks ruck evidence admissibility. \\

\textbf{DEP-1C-021} &
7.16 Operational Self-Management Systems — Tier 1 schedule stability for 6 sessions/week &
Week 2 &
Trainee &
Implement conservative substitution posture to preserve cadence; prohibit make-up doubling; \textbf{HOLD} escalation until cadence proof is restored. \\

\end{longtblr}

\endgroup

\subsection*{14. Validity Requirements (Evidence Admissibility)}
\phantomsection


Sessions:

\begin{itemize}
\item Tagged \textbf{VALID} / \textbf{MODIFIED} / \textbf{INVALID} per Stage 1.F.
\item Cardio minimum must be met continuously; failures make session \textbf{INVALID}.
\item \textbf{MODIFIED}/\textbf{INVALID} sessions require \textbf{SESSION-RC} and observable trigger notes; undocumented substitutions are treated as \textbf{INVALID}.
\end{itemize}

Tests:

\begin{itemize}
\item Tagged \textbf{VALID} / \textbf{INVALID} only.
\item Evidence Ref is required for every test event; missing Evidence Ref defaults the test to \textbf{INVALID} for gate use.
\item Environment Unit \textbf{ID} is required when applicable (\textbf{RUN} and \textbf{RUCK} tests); missing environment identity defaults \textbf{INVALID} for gate use.
\item Ruck-specific requirements for \textbf{VALID} evidence (\textbf{C12}):
  \begin{itemize}
  \item dry load verified and logged,
  \item carried water logged separately,
  \item footwear recorded,
  \item route/treadmill identity and measurement posture recorded.
  \end{itemize}
\end{itemize}

Reslot posture:

\begin{itemize}
\item Any required test that is \textbf{INVALID} is treated as not yet tested for gate decisions and triggers \textbf{HOLD}/reslot discipline per Stage 3.E.
\end{itemize}

\end{landscape}

\subsection*{15. Change Control Notes (Freeze Windows and Patch Discipline)}
\phantomsection


Allowed without patch (must be documented):

\begin{itemize}
\item Swap sequencing of tests within the same protected week to preserve spacing and safety, without adding tests.
\item Reslot an \textbf{INVALID}/missed test to the next protected window.
\item Convert a protected window into deload-only/check-in-only when readiness/environment makes gate-grade attempts inadmissible.
\item Substitute route ↔ treadmill only when Stage 2 validity/logging preserves comparability; otherwise reslot.
\end{itemize}

Requires patch:

\begin{itemize}
\item Changing gate test set (adding/removing instruments).
\item Changing protected window placement posture.
\item Extending phase duration or inserting additional protected windows.
\item Altering validity requirements or required logging fields.
\end{itemize}

\subsection*{16. Compliance Checklist (Pass and Fail)}
\phantomsection


Pass (minimum):

\begin{itemize}
\item Phase header fields complete and correct.
\item Operating constraints explicitly include:
  \begin{itemize}
  \item 6 sessions/week,
  \item required session elements in correct order,
  \item Cardio minimum 30 minutes continuous per session stated in body text,
  \item 10,000 steps/day stated in body text,
  \item admissible evidence posture (\textbf{VALID}-only gate evidence; \textbf{INVALID} provides no gate credit).
  \end{itemize}
\item Weekly structure table includes Sessions 1–6 with Cardio Modality/Intent entries containing Modality; Minutes; RPE + talk-test cue; continuity statement.
\item Protected deload/test windows stated for a 12-week phase and aligned to Stage 3 posture.
\item Test Battery references Stage 2 instruments only (C\# and \textbf{MET-2B}-\#\#\#) with no placeholders in gate-relevant rows.
\item Dependencies table includes \textbf{DEP-1C}-\#\#\# and Stage 7 categories, with owners, due-by timing, and explicit fail actions.
\item No programming content appears.
\end{itemize}

Fail (non-deployable triggers):

\begin{itemize}
\item Any gate test referenced without Stage 2 identifiers.
\item Any implication that gate-grade tests occur outside protected windows.
\item Any omission of Cardio minimum or steps baseline in operating constraints.
\item Any attempt to treat compromised/\textbf{INVALID} evidence as admissible for gate outcomes.
\end{itemize}

\end{document}